\documentclass[a4paper]{article}
\usepackage[utf8]{vntex}
\usepackage{wrapfig}
%\usepackage[english,vietnam]{babel}
%\usepackage[utf8]{inputenc}
\usepackage{mathptmx}[ptm]
%\usepackage[utf8]{inputenc}
%\usepackage[francais]{babel}
\usepackage{a4wide,amssymb,epsfig,latexsym,array,hhline,fancyhdr}
\usepackage[normalem]{ulem}
%\usepackage{soul}

\usepackage[makeroom]{cancel}
\usepackage{amsmath}
\usepackage{amsthm}
\usepackage{multicol,longtable,amscd}
\usepackage{diagbox}%Make diagonal lines in tables
\usepackage{booktabs}
\usepackage{alltt}
\usepackage[framemethod=tikz]{mdframed}% For highlighting paragraph backgrounds
\usepackage{caption,subcaption}

\usepackage{lastpage}
\usepackage[lined,boxed,commentsnumbered]{algorithm2e}
\usepackage{enumerate}
\usepackage{color}
\usepackage{graphicx}							% Standard graphics package
\graphicspath{{images/}{Images/}}
\usepackage{array}
\usepackage{tabularx, caption}
\usepackage{multirow}
\usepackage{multicol}
\usepackage{rotating}
\usepackage{graphics}
\usepackage{geometry}
\usepackage{setspace}
\usepackage{titlesec}
\usepackage{float}
\usepackage{indentfirst}
\usepackage{epsfig}
\usepackage{tikz}
\usepackage{eso-pic}

\usetikzlibrary{arrows,decorations.pathmorphing,backgrounds,calc,positioning,shapes.geometric}

% Cover frame in absolute paper coordinates so it cannot drift
\newcommand{\InternshipCoverFrame}{%
  \AddToShipoutPictureBG*{%
    \AtPageLowerLeft{%
      \begin{tikzpicture}
        \useasboundingbox (0,0) rectangle (\paperwidth,\paperheight);
        \draw[line width=4pt]
          (0.40in,0.50in) rectangle (\paperwidth-0.40in,\paperheight-0.50in);
        \draw[line width=1.5pt]
          (0.45in,0.55in) rectangle (\paperwidth-0.45in,\paperheight-0.55in);
      \end{tikzpicture}%
    }%
  }%
}

\usepackage[unicode]{hyperref}
\hypersetup{urlcolor=blue,linkcolor=black,citecolor=black,colorlinks=true} 
\usepackage{listings}

%\usepackage{pstcol} 								% PSTricks with the standard color package

\usepackage[normalem]{ulem}

\newtheorem{theorem}{{\bf Định lý}}
\newtheorem{property}{{\bf Tính chất}}
\newtheorem{proposition}{{\bf Mệnh đề}}
\newtheorem{corollary}[proposition]{{\bf Hệ quả}}
\newtheorem{lemma}[proposition]{{\bf Bổ đề}}
\theoremstyle{definition}
\newtheorem{exer}{Bài toán}

\def\thesislayout{	% A4: 210 × 297
	\geometry{
		a4paper,
		total={160mm,240mm},  % fix over page
		left=30mm,
		top=30mm,
	}
}
\def\thesisheadlayout{	% A4: 210 × 297
	\geometry{
		a4paper,
		total={160mm,240mm},  % fix over page
		left=30mm,
		top=10mm,
	}
}
\thesislayout

% Màu và kiểu hiển thị mã nguồn Python.
\definecolor{codebackground}{RGB}{247,249,252}
\definecolor{codeframe}{RGB}{203,213,225}
\definecolor{codekeyword}{RGB}{30,64,175}
\definecolor{codecomment}{RGB}{22,101,52}
\definecolor{codestring}{RGB}{180,83,9}
\definecolor{codenumber}{RGB}{100,116,139}

\lstdefinestyle{pythoncode}{
  language=Python,
  basicstyle=\small\ttfamily,
  keywordstyle=\color{codekeyword}\bfseries,
  commentstyle=\color{codecomment}\itshape,
  stringstyle=\color{codestring},
  numberstyle=\tiny\ttfamily\color{codenumber},
  numbers=left,
  stepnumber=1,
  numbersep=10pt,
  backgroundcolor=\color{codebackground},
  frame=single,
  rulecolor=\color{codeframe},
  framerule=0.6pt,
  framesep=6pt,
  xleftmargin=1.5em,
  framexleftmargin=1em,
  showspaces=false,
  showstringspaces=false,
  showtabs=false,
  keepspaces=true,
  columns=fullflexible,
  tabsize=4,
  captionpos=t,
  breaklines=true,
  breakatwhitespace=false,
  aboveskip=1em,
  belowskip=1em
}

\lstdefinestyle{terminaloutput}{
  language={},
  basicstyle=\footnotesize\ttfamily,
  backgroundcolor=\color{codebackground},
  frame=single,
  rulecolor=\color{codeframe},
  numbers=none,
  breaklines=true,
  columns=fullflexible
}

\lstset{style=pythoncode}
%\usepackage{fancyhdr}
\setlength{\headheight}{40pt}
\pagestyle{fancy}
\fancyhead{} % clear all header fields
\fancyhead[L]{
 \begin{tabular}{rl}
    \begin{picture}(25,15)(0,0)
    \put(0,-8){\includegraphics[width=10mm, height=10mm]{hcmut.png}}
    %\put(0,-8){\epsfig{width=10mm,figure=hcmut.eps}}
   \end{picture}&
	%\includegraphics[width=8mm, height=8mm]{hcmut.png} & %
	\begin{tabular}{l}
		\textbf{Trường Đại Học Bách Khoa - Đại học Quốc gia TP.HCM }\\
		\textbf{Khoa Điện - Điện Tử}
	\end{tabular} 	
 \end{tabular}
}
\fancyhead[R]{
	\begin{tabular}{l}
		\tiny \bf \\
		\tiny \bf 
	\end{tabular}  }
\fancyfoot{} % clear all footer fields
\fancyfoot[L]{\scriptsize  Báo cáo đồ án 2 - HK1 2026 - 2027}
\fancyfoot[R]{\scriptsize  Trang {\thepage}/\pageref{LastPage}}
\renewcommand{\headrulewidth}{0.3pt}
\renewcommand{\footrulewidth}{0.3pt}


%%%
\setcounter{secnumdepth}{4}
\setcounter{tocdepth}{3}
\makeatletter
\newcounter {subsubsubsection}[subsubsection]
\renewcommand\thesubsubsubsection{\thesubsubsection .\@alph\c@subsubsubsection}
\newcommand\subsubsubsection{\@startsection{subsubsubsection}{4}{\z@}%
                                     {-3.25ex\@plus -1ex \@minus -.2ex}%
                                     {1.5ex \@plus .2ex}%
                                     {\normalfont\normalsize\bfseries}}
\newcommand*\l@subsubsubsection{\@dottedtocline{3}{10.0em}{4.1em}}
\newcommand*{\subsubsubsectionmark}[1]{}
\makeatother

\everymath{\color{black}}%make in-line maths symbols blue to read/check easily

% --- Thiết lập dàn trang và dãn dòng ---
% Dãn dòng vừa phải, phù hợp với báo cáo tiếng Việt.
\setstretch{1.15}
\raggedbottom
\setlength{\parindent}{1.25em}
\setlength{\parskip}{0.35em plus 0.08em minus 0.04em}
\setlength{\emergencystretch}{2em}

% Khoảng cách tiêu đề nhất quán, tránh tiêu đề bị dính vào đoạn văn.
\titleformat{\section}
  {\normalfont\Large\bfseries}
  {\thesection}{0.75em}{}
\titlespacing*{\section}{0pt}{2.2ex plus 0.6ex minus 0.2ex}{1.1ex plus 0.2ex}

\titleformat{\subsection}
  {\normalfont\large\bfseries}
  {\thesubsection}{0.75em}{}
\titlespacing*{\subsection}{0pt}{1.8ex plus 0.5ex minus 0.2ex}{0.8ex plus 0.2ex}

\titleformat{\subsubsection}
  {\normalfont\normalsize\bfseries}
  {\thesubsubsection}{0.75em}{}
\titlespacing*{\subsubsection}{0pt}{1.4ex plus 0.4ex minus 0.2ex}{0.5ex plus 0.1ex}

% Khoảng cách hình, bảng và chú thích cân đối hơn.
\captionsetup[figure]{labelfont={small,bf},textfont={small,it},skip=5pt}
\setlength{\floatsep}{10pt plus 2pt minus 2pt}
\setlength{\textfloatsep}{12pt plus 2pt minus 3pt}
\setlength{\intextsep}{10pt plus 2pt minus 2pt}

\thesislayout

\begin{document}

\hypersetup{pageanchor=false}
\newgeometry{left=25mm,right=25mm,top=20mm,bottom=18mm}
\begin{titlepage}
\InternshipCoverFrame
\thispagestyle{empty}
\begingroup
\centering
\setstretch{1}
\setlength{\parskip}{0pt}
\bfseries
\fontsize{13}{18}\selectfont
TRƯỜNG ĐẠI HỌC BÁCH KHOA TP. HỒ CHÍ MINH\par
\vspace{4mm}
KHOA ĐIỆN -- ĐIỆN TỬ\par
\vspace{4mm}
BỘ MÔN ĐIỀU KHIỂN TỰ ĐỘNG\par

\vspace{17mm}
\includegraphics[width=40mm]{hcmut.png}\par
\vspace{18mm}
NGUYỄN VĂN TUẤN ANH -- 2310131

\vspace{28mm}
{\fontsize{15}{20}\selectfont ĐỒ ÁN 2\par}
\vspace{5mm}
{\fontsize{16}{24}\selectfont
THIẾT KẾ  UAV 4 CÁNH TỰ ĐỘNG HẠ CÁNH VÀ GẮP HÀNG BẰNG MÃ ARUCO \par}
\vspace{6mm}
DESIGN DRONE AUTO LANDING AND PICKUP USING ARUCO MARKERS \par

\vspace{18mm}
{\fontsize{11}{16}\selectfont
GIẢNG VIÊN HƯỚNG DẪN\par
\vspace{4mm}
PGS. TS NGUYỄN TẤN LŨY\par}

\vfill
{\fontsize{12}{16}\selectfont TP. HỒ CHÍ MINH, 2026\par}
\endgroup
\end{titlepage}
\restoregeometry
% Trang đề cương chi tiết theo mẫu của Bộ môn Điều khiển tự động.
\clearpage
\newgeometry{left=25mm,right=15mm,top=25mm,bottom=20mm}
\thispagestyle{empty}
\begingroup
\setstretch{1}
\setlength{\parindent}{0pt}
\setlength{\parskip}{0pt}
\fontsize{12}{16}\selectfont

\noindent
\begin{minipage}[t]{0.51\textwidth}
  \centering
  {\fontsize{10}{14}\selectfont
  TRƯỜNG ĐẠI HỌC BÁCH KHOA TP. HỒ CHÍ MINH\par}
  \vspace{1mm}
  \textbf{KHOA ĐIỆN -- ĐIỆN TỬ}\par
  \textbf{BỘ MÔN: ĐIỀU KHIỂN TỰ ĐỘNG}
\end{minipage}\hfill
\begin{minipage}[t]{0.47\textwidth}
  \centering
  {\fontsize{10}{14}\selectfont
  CỘNG HÒA XÃ HỘI CHỦ NGHĨA VIỆT NAM\par
  \vspace{1mm}
  Độc lập - Tự do - Hạnh phúc\par
  \rule{0.6\linewidth}{0.4pt}\par} 
  \vspace{4mm}
  {\fontsize{10}{14}\selectfont\itshape
  TP. HCM, ngày....tháng....năm....}
\end{minipage}

\vspace{8mm}
\begin{center}
  \textbf{ĐỀ CƯƠNG CHI TIẾT}
\end{center}
\vspace{1mm}

% Các chiều cao cố định giữ nguyên khoảng trống để điền và ký xác nhận.
\setlength{\tabcolsep}{2.5mm}
\renewcommand{\arraystretch}{1}
\noindent
\begin{tabular}{|p{\dimexpr\textwidth-2\tabcolsep-2\arrayrulewidth\relax}|}
\hline
\begin{minipage}[c][11mm][c]{\linewidth}
  \textbf{TÊN ĐỒ ÁN 2:}
\end{minipage}\\ \hline
\begin{minipage}[c][10mm][c]{\linewidth}
  \textbf{Cán bộ hướng dẫn: PGS. TS NGUYỄN TẤN LŨY}
\end{minipage}\\ \hline
\begin{minipage}[c][12mm][c]{\linewidth}
  \textbf{Thời gian thực hiện:} Từ ngày....................đến ngày................
\end{minipage}\\ \hline
\begin{minipage}[c][23mm][c]{\linewidth}
  \textbf{Sinh viên thực hiện:}\par
  \vspace{4mm}
  \textbf{NGUYỄN VĂN TUẤN ANH -- 2310131}
\end{minipage}\\ \hline
\begin{minipage}[c][76mm][t]{\linewidth}
  \vspace{3mm}
  \textbf{Nội dung đề tài:} \textit{(Mô tả chi tiết}\par
  \vspace{5mm}
  \textit{mục tiêu,}\par
  \vspace{5mm}
  \textit{phạm vi, đối tượng,}\par
  \vspace{5mm}
  \textit{phương pháp thực hiện,}\par
  \vspace{5mm}
  \textit{kết quả mong đợi của đề tài)}
\end{minipage}\\ \hline
\begin{minipage}[c][42mm][t]{\linewidth}
  \vspace{3mm}
  \textbf{Kế hoạch thực hiện:} \textit{(Mô tả kế hoạch làm việc và phân công công việc cho từng sinh viên tham gia)}
\end{minipage}\\ \hline
\multicolumn{1}{|@{}c@{}|}{%
  \begin{minipage}[c][25mm][t]{\dimexpr0.5\textwidth-\arrayrulewidth\relax}
    \vspace{4mm}\centering
    \textbf{Xác nhận của Cán bộ hướng dẫn}\\
    (ký và ghi rõ họ tên)
  \end{minipage}%
  \vrule width\arrayrulewidth
  \begin{minipage}[c][25mm][t]{\dimexpr0.5\textwidth-2\arrayrulewidth\relax}
    \vspace{4mm}\centering
    {\fontsize{11}{15}\selectfont TP. HCM, ngày....tháng ....năm.....}\\
    Sinh viên\\
    (ký và ghi rõ họ tên)
  \end{minipage}%
}\\ \hline
\end{tabular}
\par
\endgroup
\clearpage
\restoregeometry

\hypersetup{pageanchor=true}
\pagenumbering{arabic}

\tableofcontents

\newpage
\listoffigures

\newpage
\listoftables

\newpage
% Trang đề cương chi tiết theo mẫu của Bộ môn Điều khiển tự động.
\clearpage
\newgeometry{left=25mm,right=15mm,top=25mm,bottom=20mm}
\thispagestyle{empty}
\begingroup
\setstretch{1}
\setlength{\parindent}{0pt}
\setlength{\parskip}{0pt}
\fontsize{12}{16}\selectfont

\noindent
\begin{minipage}[t]{0.51\textwidth}
  \centering
  {\fontsize{10}{14}\selectfont
  TRƯỜNG ĐẠI HỌC BÁCH KHOA TP. HỒ CHÍ MINH\par}
  \vspace{1mm}
  \textbf{KHOA ĐIỆN -- ĐIỆN TỬ}\par
  \textbf{BỘ MÔN: ĐIỀU KHIỂN TỰ ĐỘNG}
\end{minipage}\hfill
\begin{minipage}[t]{0.47\textwidth}
  \centering
  {\fontsize{10}{14}\selectfont
  CỘNG HÒA XÃ HỘI CHỦ NGHĨA VIỆT NAM\par
  \vspace{1mm}
  Độc lập - Tự do - Hạnh phúc\par
  \rule{0.6\linewidth}{0.4pt}\par} 
  \vspace{4mm}
  {\fontsize{10}{14}\selectfont\itshape
  TP. HCM, ngày....tháng....năm....}
\end{minipage}

\vspace{8mm}
\begin{center}
  \textbf{ĐỀ CƯƠNG CHI TIẾT}
\end{center}
\vspace{1mm}

% Các chiều cao cố định giữ nguyên khoảng trống để điền và ký xác nhận.
\setlength{\tabcolsep}{2.5mm}
\renewcommand{\arraystretch}{1}
\noindent
\begin{tabular}{|p{\dimexpr\textwidth-2\tabcolsep-2\arrayrulewidth\relax}|}
\hline
\begin{minipage}[c][11mm][c]{\linewidth}
  \textbf{TÊN ĐỒ ÁN 2:}
\end{minipage}\\ \hline
\begin{minipage}[c][10mm][c]{\linewidth}
  \textbf{Cán bộ hướng dẫn: PGS. TS NGUYỄN TẤN LŨY}
\end{minipage}\\ \hline
\begin{minipage}[c][12mm][c]{\linewidth}
  \textbf{Thời gian thực hiện:} Từ ngày....................đến ngày................
\end{minipage}\\ \hline
\begin{minipage}[c][23mm][c]{\linewidth}
  \textbf{Sinh viên thực hiện:}\par
  \vspace{4mm}
  \textbf{NGUYỄN VĂN TUẤN ANH -- 2310131}
\end{minipage}\\ \hline
\begin{minipage}[c][76mm][t]{\linewidth}
  \vspace{3mm}
  \textbf{Nội dung đề tài:} \textit{(Mô tả chi tiết}\par
  \vspace{5mm}
  \textit{mục tiêu,}\par
  \vspace{5mm}
  \textit{phạm vi, đối tượng,}\par
  \vspace{5mm}
  \textit{phương pháp thực hiện,}\par
  \vspace{5mm}
  \textit{kết quả mong đợi của đề tài)}
\end{minipage}\\ \hline
\begin{minipage}[c][42mm][t]{\linewidth}
  \vspace{3mm}
  \textbf{Kế hoạch thực hiện:} \textit{(Mô tả kế hoạch làm việc và phân công công việc cho từng sinh viên tham gia)}
\end{minipage}\\ \hline
\multicolumn{1}{|@{}c@{}|}{%
  \begin{minipage}[c][25mm][t]{\dimexpr0.5\textwidth-\arrayrulewidth\relax}
    \vspace{4mm}\centering
    \textbf{Xác nhận của Cán bộ hướng dẫn}\\
    (ký và ghi rõ họ tên)
  \end{minipage}%
  \vrule width\arrayrulewidth
  \begin{minipage}[c][25mm][t]{\dimexpr0.5\textwidth-2\arrayrulewidth\relax}
    \vspace{4mm}\centering
    {\fontsize{11}{15}\selectfont TP. HCM, ngày....tháng ....năm.....}\\
    Sinh viên\\
    (ký và ghi rõ họ tên)
  \end{minipage}%
}\\ \hline
\end{tabular}
\par
\endgroup
\clearpage
\restoregeometry

\newpage
\section{Tóm tắt lý thuyết}


\newpage
\section{Example 1}

\newpage
\section{Example 2 - Image Compression}

\end{document}
